\section{问题定义}
\label{sec:problem}
% 输入对象、决策、策略、优化问题。每个定义后跟贯穿例子。

\begin{definition}[候选集]
\label{def:candidate}
给定表 $T$，候选集 $\cand$ 是<定义>。
\end{definition}

\begin{example}
\label{ex:candidate}
在例~\ref{ex:motivation}中，<贯穿例子的候选集>。
\end{example}

\begin{definition}[标注策略]
\label{def:policy}
策略 $\policy$ 把<状态>映射为<决策>，且至多标注 $\budget$ 个候选。
\end{definition}

目标为
\begin{equation}
\label{eq:objective}
\max_{\policy}\ \mathrm{Quality}(\policy(\cand)) \quad \text{s.t.} \quad |\policy(\cand)| \le \budget .
\end{equation}

\begin{theorem}[困难性]
\label{thm:hard}
在一般的逐候选代价下，问题~\eqref{eq:objective}的判定版本是 NP 完全的。
\end{theorem}
证明见附录~\ref{app:proof-hard}。这一结果说明第~\ref{sec:method}~节需要近似方法。
